\BourbakiBackSection{BIBLIOGRAPHY}

{[}1{]} B. BOLZANO, Rein analytischer Beweis des Lehrsatzes, dass zwischen je zwei Werthen, die ein entgegengesetztes Resultat gewähren, wenigstens eine reelle Wurzel liegt, Ostwald's Klassiker, no. 153, Leipzig, 1905.

{[}2{]} A.-L. CAUCHY, Sur la convergence des séries {[}Exercices d'Analyse, 2 \(^{e}\) Année, Paris, 1827, p.~221 = Oeuvres (II), vol.~7, Paris (Gauthier-Villars) 1889, p.~267{]}.

{[}3{]} G. CANTOR, Gesammelte Abhandlungen, Berlin (Springer), 1932.

{[}4{]} E. HEINE, Über trigonometrische Reihen, Crelle's Journal, 71 (1870), pp.~353-365.

{[}5{]} E. HEINE, Die Elemente der Functionenlehre, Crelle's Journal, 74 (1872), pp.~172-188.

{[}6{]} M. FRÉCHET, Sur quelques points du calcul fonctionnel, Rend. Palermo, 22 (1906), pp.~1-74.

{[}7{]} M. FRÉCHET, Les ensembles abstraits et le calcul fonctionnel, Rend. Palermo, 30 (1910), pp.~1-26.

{[}8{]} F. HAUSDORFF, Grundzüge der Mengenlehre, Leipzig (Veit), 1914.

{[}8 a{]} F. HAUSDORFF, Mengenlehre, Berlin (de Gruyter), 1927.

{[}9{]} A. WEIL, Sur les espaces à structure uniforme et sur la topologie générale, Act. Scient. et Ind., no. 551, Paris (Hermann), 1937.

{[}10{]} P. ALEXANDROFF and H. HOPF, Topologie I, Berlin (Springer), 1935.

